\chapter{师范生数学教学技能调查研究与问卷编制}


\section{研究对象}

本研究以电子问卷调查的方式，利用问卷星制作问卷，以在线方式进行调查。本研究以华中师范大学2019、2018两届数学师范生为样本，问卷共收集185份，有效问卷177份，有效率95.67\%。受试者的基本人口学变量资料列于表~\ref{table:被试基本人口学变量信息} 。

\setcounter{table}{0}
\begin{table}[htbp]
  \centering
  \caption{被试基本人口学变量信息}
  \label{table:被试基本人口学变量信息}
  \begin{tblr}{
    width = 0.7\textwidth,
    colspec = {X[2,c]X[2,c]X[1,c]X[1,c]},
    hline{1,2,Z} = {solid},
    cell{2}{1} = {r=2}{m},
    cell{4}{1} = {r=2}{m},
    cell{6}{1} = {r=3}{m},
    cell{9}{1} = {r=4}{m}
  }
    变量    & 选项    & 频率  & 百分比  \\
    性别    & 女     & 141 & 80\% \\
          & 男     & 36  & 20\% \\
    年级    & 大四    & 106 & 60\% \\
          & 大三    & 71  & 40\% \\
    计算机等级 & 一级    & 72  & 41\% \\
          & 二级    & 95  & 54\% \\
          & 三级及以上 & 10  & 6\%  \\
    普通话等级 & 一甲    & 2   & 1\%  \\
          & 一乙    & 43  & 24\% \\
          & 二甲    & 90  & 51\% \\
          & 二乙    & 42  & 24\%
  \end{tblr}
\end{table}

\section{研究工具}

\subsection{问卷的编制}

本研究是在对相关文献、书籍和微格教学技能训练手册进行分析、归纳、总结的基础上，形成问卷。同时，还与数学教育方面的专家、数学师范生进行了沟通，排除了问卷中的个别问题，并对部分问题的表述进行了修正，形成了最终的调查问卷（详见附录）。

问卷由三个部分组成。第一个部分为数学师范生基本信息，包括性别、普通话等级、大学英语等级、计算机等级；第二部分从五个不同的维度对数学专业师范生的数学教学技能的掌握状况进行了调查，一共39道题。第三部分为师范生数学教学技能的训练情况调查，共10道题。调查问卷的主体部分采用5点记分法，“1→5”表示“不具备→完全具备”，打分越高，说明教学技能掌握情况越好。


\subsection{问卷信度}

正式的问卷量表调查数据输入SPSS26.0进行信度分析。问卷整体的基于标准化项的克隆巴赫Alpha系数为0.935，达到了很好的信度水平。其中各项教学技能的克隆巴赫Alpha系数分布如表~\ref{table:问卷信度分布表} 。

\begin{table}[htbp]
  \centering
  \caption{问卷信度分布表}
  \label{table:问卷信度分布表}
  \begin{tblr}{
    width = 0.7\textwidth,
    colspec = {X[2,c]X[2,c]X[1,c]},
    hline{1,2,Z} = {solid},
  }
    项目     & 克隆巴赫Alpha系数 & 项目数 \\
    教学设计技能 & 0.796       & 5   \\
    数学说课技能 & 0.756       & 4   \\
    课堂教学技能 & 0.911       & 24  \\
    教学评价技能 & 0.746       & 3   \\
    教学研究技能 & 0.725       & 3   \\
    问卷量表   & 0.935       & 39 
  \end{tblr}
\end{table}


表~\ref{table:问卷信度分布表} 表明，在各项教学技能上，每个分量表的克隆巴赫Alpha系数均高于0.7，因此问卷信度可以接受。

\begin{table}[htbp]
  \centering
  \caption{课堂教学技能量表信度分布表}
  \label{table:课堂教学技能量表信度分布表}
  \begin{tblr}{
    width = 0.7\textwidth,
    colspec = {X[2,c]X[2,c]X[1,c]},
    hline{1,2,Z} = {solid},
  }
    项目     & 克隆巴赫Alpha系数 & 项目数 \\
    语言技能   & 0.677       & 3   \\
    板书技能   & 0.792       & 3   \\
    讲解技能   & 0.751       & 4   \\
    提问技能   & 0.611       & 3   \\
    导入技能   & 0.704       & 4   \\
    结束技能   & 0.652       & 3   \\
    媒体运用技能 & 0.669       & 4   \\
  \end{tblr}
\end{table}


表~\ref{table:课堂教学技能量表信度分布表} 表明，在各项课堂教学技能上，每个分量表的克隆巴赫Alpha系数均高于0.6，因此课堂教学技能量表信度可以接受。


\subsection{问卷效度}

本问卷效度通过KMO取样适切性量数和巴特利特球形检验值来检验。

\begin{table}[htbp]
  \centering
  \caption{KMO和巴特利特检验}
  \label{table:KMO和巴特利特检验}
  \begin{tblr}{
    width = 0.7\textwidth,
    colspec = {X[2,c]X[2,c]X[1,c]},
    hline{1,2,Z} = {solid},
    cell{1}{1} = {c=2}{c},
    cell{2}{1} = {r=3}{m}
  }
    KMO取样适切性量数 &      & 0.879    \\
    巴特利特球形度检验   & 近似卡方 & 3057.278 \\
                & 自由度  & 741      \\
                & 显著性  & 0.000   
  \end{tblr}
\end{table}


表~\ref{table:KMO和巴特利特检验} 表明，KMO检验的系数结果为0.879，大于0.5，且巴特利特球形检验的显著性是0.000，说明本问卷的效度较好。


\subsection{问卷的数据处理}

本研究使用SPSS26.0软件对数据进行分析.